% Part: first-order-logic
% Chapter: introduction
% Section: first-order-logic

\documentclass[../../../include/open-logic-section]{subfiles}

\begin{document}

\olfileid{fol}{int}{fol}

\olsection{Eerste-ordelogica}

Waarschijnlijk ken je de eerste-ordelogica al uit een eerste inleiding
tot de formele logica.\footnote{Daar gaan we eigenlijk min of meer van
uit!{} Is dat niet zo, dan kun je eerst een eenvoudiger leerboek
bekijken, zoals \emph{forall x} \citep{Magnus2021}.} Misschien ken je
haar als ``kwantorlogica'' of ``predicatenlogica''. De eerste-ordelogica
is om te beginnen een formele taal. Zo'n taal heeft een vaste
woordenschat. Haar uitdrukkingen zijn reeksen symbolen uit die
woordenschat, maar niet iedere reeks is toegestaan. Er zijn verschillende
soorten toegestane uitdrukkingen: termen, !!{formula}s en !!{sentence}s.
Het gaat ons vooral om !!{sentence}s van de eerste-ordelogica. Ze zijn de
formele tegenhangers van gewone Nederlandse zinnen. We kunnen er de
vragen over stellen die voor een logicus belangrijk zijn. Bijvoorbeeld:
\begin{itemize}
    \item Volgt $!B$ logisch uit~$!A$?
    \item Is $!A$ logisch waar, logisch onwaar of
contingent?
    \item Zijn $!A$ en $!B$ gelijkwaardig?
\end{itemize}

Deze vragen gaan vooral over de ``betekenis'' van !!{sentence}s in de
eerste-ordelogica. Neem de vraag of $!B$ logisch uit~$!A$ volgt. Een
filosoof vraagt dan: bestaat er een geval waarin $!A$ waar is maar $!B$
onwaar ($!B$ volgt dan niet uit~$!A$)? Of maakt ieder geval dat $!A$ waar
maakt ook $!B$ waar ($!B$ volgt dan wel uit~$!A$)? We weten alleen nog
niet wat zo'n ``geval'' precies is. Daarvoor dient de \emph{semantiek}:
de betekenisregels. De semantiek van de eerste-ordelogica maakt het idee
van een ``geval'' wiskundig precies. Ook maakt zij precies wat het
betekent dat !!a{sentence}~$!A$ \emph{waar is in} zo'n geval. Het
wiskundige model van een geval dat we gaan maken, noemen we
!!a{structure}. De relatie die ``waar in'' precies maakt, heet
\emph{vervulling}. We gaan dus vastleggen wat ``$!A$ wordt vervuld
in~$\Struct{M}$'' betekent (in symbolen: $\Sat{M}{!A}$), voor
!!{sentence}s~$!A$ en !!{structure}s~$\Struct{M}$. Daarna kunnen we ook
``volgt uit'' en ``is logisch waar'' precies definiëren. Met die
definities kunnen we
wiskundig precies bepalen of bijvoorbeeld
$\lforall[x][(!A(x) \lif !B(x)),
\lexists[x][!A(x)] \Entails \lexists[x][!B(x)]]$. Het antwoord is
natuurlijk ``ja''. Als je al hebt geleerd gewone zinnen in
eerste-ordelogica te zetten, herken je hierin bijvoorbeeld: ``Alle
mieren zijn insecten; er zijn mieren; dus zijn er insecten.'' Dat is
duidelijk een geldig argument. Ons wiskundige model van ``volgt uit''
voor de formele taal moet daarom hetzelfde antwoord geven.

Uit een eerste cursus formele logica ken je waarschijnlijk ook
\emph{!!{derivation}s}. Daar heb je vermoedelijk geoefend om zulke
!!{derivation}s te vinden. Misschien heb je zelfs !!a{derivation} gemaakt
voor de gevolgtrekking hierboven. Een !!{derivation} kan op veel manieren
worden gegeven. Je hebt misschien gewerkt met ``natuurlijke deductie''
of ``semantische tableaus'', maar er bestaan nog veel andere systemen.
Met systemen voor formele !!{derivation} kunnen we de vragen hierboven beantwoorden.
Neem bijvoorbeeld !!a{derivation} in natuurlijke deductie waarin
$\lforall[x][(!A(x) \lif !B(x))$ en $\lexists[x][!A(x)]$ de premissen
zijn en $\lexists[x][!B(x)]]$ de conclusie, dus de laatste regel, is.
Deze afleiding \emph{controleert} dat $\lexists[x][!B(x)]$ logisch volgt
uit $\lforall[x][(!A(x) \lif !B(x))]$ en $\lexists[x][!A(x)]$.

Maar waarom bewijst zo'n afleiding dat? Op het eerste gezicht hebben
systemen voor formele !!{derivation} niets met semantiek te maken. Bij een formele
!!{derivation} zet je symbolen volgens bepaalde regels neer. Over
``gevallen'' of ``waar in'' gaat dat niet. We moeten dus apart, op het
niveau van de logica zelf, bewijzen hoe systemen voor formele !!{derivation} en
semantiek samenhangen. We hebben bijvoorbeeld een wiskundig bewijs nodig
van het volgende. Een formele !!{derivation} van $\lexists[x][!B(x)]$
uit de premissen $\lforall[x][(!A(x) \lif !B(x))]$ en
$\lexists[x][!A(x)]$ is mogelijk precies als
$\lforall[x][(!A(x) \lif !B(x))$ en $\lexists[x][!A(x)]$ samen
$\lexists[x][!B(x)]]$ impliceren. Voordat we dat kunnen bewijzen, moeten
we eerst veel precies werk doen: de syntaxis, dus de vormregels, en de
semantiek, dus de betekenisregels, moeten goed zijn gedefinieerd.

\end{document}
